\section{Experiments}
We separate two questions: does training on expanded graphs improve the simulator, and does residual risk identify a better placement of a fixed expansion budget? Fresh Goop training and paired WaterDrop continuations test both. Earlier WaterDrop controls examine why error ranking can disagree with action value. The original Sand and WaterDrop comparisons remain in Appendix~\ref{sec:historical-audit}.

\subsection{Learning to use expanded graphs}
\begin{figure*}[t]
\centering
\setlength{\unitlength}{1pt}
\begin{picture}(470,213)
\put(30.000000,194.000000){\makebox(0,0)[l]{\small\bfseries A  Graph exposure}}
\put(268.000000,194.000000){\makebox(0,0)[l]{\small\bfseries B  Risk placement}}
\put(30.000000,179.000000){\makebox(0,0)[l]{\footnotesize Full-rollout MSE under random25}}
\put(268.000000,179.000000){\makebox(0,0)[l]{\footnotesize Observed-state risk minus random MSE}}
\put(268.000000,166.000000){\makebox(0,0)[l]{\scriptsize Position-coordinate MSE ($\times 10^{-10}$)}}
\put(30.000000,166.000000){\makebox(0,0)[l]{\scriptsize Mean over all 395 forecasts}}
\linethickness{0.35pt}
\qbezier(43.000000,51.000000)(43.000000,103.000000)(43.000000,155.000000)
\linethickness{0.55pt}
\qbezier(43.000000,51.000000)(133.000000,51.000000)(223.000000,51.000000)
\linethickness{0.35pt}
\qbezier(40.000000,51.000000)(41.500000,51.000000)(43.000000,51.000000)
\put(36.000000,51.000000){\makebox(0,0)[r]{\scriptsize 0}}
\linethickness{0.35pt}
\qbezier(40.000000,88.818182)(41.500000,88.818182)(43.000000,88.818182)
\put(36.000000,88.818182){\makebox(0,0)[r]{\scriptsize 0.2}}
\linethickness{0.35pt}
\qbezier(40.000000,126.636364)(41.500000,126.636364)(43.000000,126.636364)
\put(36.000000,126.636364){\makebox(0,0)[r]{\scriptsize 0.4}}
\linethickness{0.35pt}
\qbezier(81.000000,89.381448)(129.500000,79.038378)(178.000000,68.695309)
\linethickness{0.5pt}
\put(81.000000,89.381448){\circle{4.4}}
\linethickness{0.5pt}
\put(178.000000,68.695309){\circle{4.4}}
\linethickness{0.35pt}
\qbezier(85.000000,89.464997)(133.500000,84.725765)(182.000000,79.986532)
\linethickness{0.5pt}
\put(85.000000,89.464997){\circle*{4.4}}
\linethickness{0.5pt}
\put(182.000000,79.986532){\circle*{4.4}}
\linethickness{0.35pt}
\qbezier[20](89.000000,144.933322)(137.500000,111.119071)(186.000000,77.304820)
\linethickness{0.5pt}
\put(89.000000,144.933322){\makebox(0,0){\scriptsize$\triangle$}}
\linethickness{0.5pt}
\put(186.000000,77.304820){\makebox(0,0){\scriptsize$\triangle$}}
\linethickness{1.20pt}
\qbezier(77.000000,107.926589)(85.000000,107.926589)(93.000000,107.926589)
\linethickness{1.20pt}
\qbezier(174.000000,75.328887)(182.000000,75.328887)(190.000000,75.328887)
\put(85.000000,37.000000){\makebox(0,0){\scriptsize Base-only}}
\put(85.000000,27.000000){\makebox(0,0){\scriptsize training}}
\put(182.000000,37.000000){\makebox(0,0){\scriptsize Mixed-graph}}
\put(182.000000,27.000000){\makebox(0,0){\scriptsize training}}
\linethickness{0.35pt}
\qbezier(281.000000,51.000000)(281.000000,103.000000)(281.000000,155.000000)
\linethickness{0.55pt}
\qbezier(281.000000,51.000000)(371.000000,51.000000)(461.000000,51.000000)
\linethickness{0.35pt}
\qbezier(278.000000,51.000000)(279.500000,51.000000)(281.000000,51.000000)
\put(274.000000,51.000000){\makebox(0,0)[r]{\scriptsize 0}}
\linethickness{0.35pt}
\qbezier(278.000000,85.666667)(279.500000,85.666667)(281.000000,85.666667)
\put(274.000000,85.666667){\makebox(0,0)[r]{\scriptsize 3}}
\linethickness{0.35pt}
\qbezier(278.000000,120.333333)(279.500000,120.333333)(281.000000,120.333333)
\put(274.000000,120.333333){\makebox(0,0)[r]{\scriptsize 6}}
\linethickness{0.35pt}
\qbezier(278.000000,155.000000)(279.500000,155.000000)(281.000000,155.000000)
\put(274.000000,155.000000){\makebox(0,0)[r]{\scriptsize 9}}
\linethickness{0.35pt}
\qbezier(319.000000,112.461269)(367.500000,84.862955)(416.000000,57.264641)
\linethickness{0.5pt}
\put(319.000000,112.461269){\circle{4.4}}
\linethickness{0.5pt}
\put(416.000000,57.264641){\circle{4.4}}
\linethickness{0.35pt}
\qbezier(323.000000,143.936119)(371.500000,137.970393)(420.000000,132.004666)
\linethickness{0.5pt}
\put(323.000000,143.936119){\circle*{4.4}}
\linethickness{0.5pt}
\put(420.000000,132.004666){\circle*{4.4}}
\linethickness{0.35pt}
\qbezier[20](327.000000,118.034885)(375.500000,90.034515)(424.000000,62.034144)
\linethickness{0.5pt}
\put(327.000000,118.034885){\makebox(0,0){\scriptsize$\triangle$}}
\linethickness{0.5pt}
\put(424.000000,62.034144){\makebox(0,0){\scriptsize$\triangle$}}
\linethickness{1.20pt}
\qbezier(315.000000,124.810758)(323.000000,124.810758)(331.000000,124.810758)
\linethickness{1.20pt}
\qbezier(412.000000,83.767817)(420.000000,83.767817)(428.000000,83.767817)
\put(323.000000,37.000000){\makebox(0,0){\scriptsize Base-only}}
\put(323.000000,27.000000){\makebox(0,0){\scriptsize training}}
\put(420.000000,37.000000){\makebox(0,0){\scriptsize Mixed-graph}}
\put(420.000000,27.000000){\makebox(0,0){\scriptsize training}}
\linethickness{0.5pt}
\put(116.000000,10.000000){\circle{4.4}}
\put(124.000000,10.000000){\makebox(0,0)[l]{\scriptsize Seed 0}}
\linethickness{0.5pt}
\put(188.000000,10.000000){\circle*{4.4}}
\put(196.000000,10.000000){\makebox(0,0)[l]{\scriptsize Seed 1}}
\linethickness{0.5pt}
\put(260.000000,10.000000){\makebox(0,0){\scriptsize$\triangle$}}
\put(268.000000,10.000000){\makebox(0,0)[l]{\scriptsize Seed 2}}
\linethickness{1.20pt}
\qbezier(325.000000,10.000000)(332.000000,10.000000)(339.000000,10.000000)
\put(344.000000,10.000000){\makebox(0,0)[l]{\scriptsize Mean}}
\end{picture}
\caption{Goop separates graph exposure from risk placement. Three paired faithful models per training arm are trained from initialization for 100,000 updates. A: random25 evaluation improves in every seed after mixed-graph training (30 complete H395 rollouts per seed and arm). B: previous-observed-base risk remains worse than random at identical histories in every seed, although its deficit narrows (150 histories per model). Symbols show seed means; short bars show three-seed means; small horizontal offsets separate seeds. These panels use different endpoints and scales. The full study records 1,077/1,080 completed rollouts: three seed-2 mixed-arm guards affect base, dense and cached risk. The mixed cached-risk full-horizon mean and its training interaction remain undefined.}
\label{fig:goop-exposure-placement}
\end{figure*}

We train six Goop models from initialization for 100,000 updates: three paired seeds each for base-only and mixed-graph training. All use faithful regression, width 128, ten message-passing blocks and all 1,000 official training trajectories. The mixed arm adds a uniformly random quarter of the optional pairs independently with probability $1/2$. Initial weights, sampled frames and noise are paired by seed; the native capped base graph and self-messages are preserved. At evaluation, random25, speed25, risk25 and relative-velocity-RMS25 use the same quarter-annulus budget, with radius ratio 1.267. Base and dense control the amount of expansion. Observed-state tests use 150 common histories per model, scoring risk from the preceding observed history under the base graph. Autonomous rollouts feed predictions back as inputs; cached risk carries scores from its own preceding selected graph.

Graph exposure improves full-rollout error under an unchanged random25 evaluation policy: mean H395 position MSE falls from $0.301\pm0.169$ to $0.129\pm0.031$, and every paired seed improves. The matched-history comparison tests exposure and placement before autonomous feedback (Table~\ref{tab:goop-exposure-placement}). Previous-observed risk remains worse than random in every mixed-model seed, despite a smaller deficit than with base-only training. Its mean frame Spearman correlation is $.214\pm.088$ with base error but $.001\pm.048$ with its own sparse-action benefit. Exposure therefore improves this simulator without making residual rank a reliable guide to useful extra messages.

\begin{table}[t]
\centering\scriptsize
\caption{Observed-test position-coordinate MSE contrasts. Each column uses its stated coordinate-squared scale. Entries are means $\pm$ sample SDs of three paired seeds. Negative contrasts favor the first term; the last row changes the risk-minus-random gap. Goop and Sand train from initialization to 100k; WaterDrop continues inspected 100k parents to 110k. Materials are not pooled.}
\label{tab:goop-exposure-placement}
\setlength{\tabcolsep}{2pt}
\begin{tabular}{lccc}
\toprule
Contrast & \shortstack{Goop\\$\times10^{-10}$} & \shortstack{WaterDrop\\$\times10^{-10}$} & \shortstack{Sand\\$\times10^{-9}$}\\
\midrule
Mix $-$ base (base) & $-3.04\pm5.12$ & $+0.70\pm0.40$ & $-0.84\pm2.91$\\
Mix $-$ base (random) & $-5.51\pm4.28$ & $+0.28\pm0.92$ & $-3.85\pm3.66$\\
Risk $-$ random (base) & $+6.39\pm1.45$ & $+1.62\pm1.07$ & $+4.76\pm3.80$\\
Risk $-$ random (mix) & $+2.84\pm3.62$ & $-0.99\pm0.67$ & $+1.87\pm1.04$\\
Training interaction & $-3.55\pm2.18$ & $-2.61\pm0.74$ & $-2.89\pm2.81$\\
\bottomrule
\end{tabular}
\end{table}

We average frames within trajectories, then trajectories within each seed; all three seed values enter the reported means and sample SDs. Figure~\ref{fig:goop-exposure-placement} shows every paired seed. The Goop study follows the WaterDrop analyses and is exploratory; its endpoints and policies were fixed before evaluation.

A separate WaterDrop follow-up pairs 10k base-only and mixed-graph continuations from three faithful 100k parents, inheriting optimizer state and matched frame/noise schedules. Its exposure effect is conditional on those parents, distinct from Goop training from initialization. Mixed training changes the observed-test risk-minus-random position-MSE gap from $(+1.62\pm1.07)\times10^{-10}$ to $(-0.99\pm0.67)\times10^{-10}$, with every seed changing sign (Table~\ref{tab:goop-exposure-placement}). Validation's relative gap also improves, but mixed risk remains worse than random on average.

\subsection{Does the effect survive autonomous feedback?}
All six policies are evaluated for all 395 forecasts on 30 Goop test trajectories per model. Of 1,080 required outcomes, 1,077 complete and three hit the candidate-pair guard, all in mixed seed 2: base, dense and cached risk each fail once. Their three-seed full-horizon means remain undefined (Table~\ref{tab:goop-rollouts}), as does the full-horizon risk-minus-random training interaction. The cached-risk comparison with random changes sign across the two defined seeds; we do not average those survivors. Pointwise MSE at forecast 200 remains a secondary endpoint and does not replace the full horizon.

\begin{table}[t]
\centering\scriptsize
\caption{Full-horizon position-MSE training effects (mixed minus base; three-seed mean $\pm$ sample SD). Goop H395 and Sand H314 require 90 outcomes per arm/policy; WaterDrop H995 requires 81. Each lower line reports base/mixed unsuccessful ($U$) and not-completed ($N$) counts. A dash preserves an undefined full-horizon effect; NP denotes a policy not predeclared.}
\label{tab:goop-rollouts}
\setlength{\tabcolsep}{2pt}
\begin{tabular}{lccc}
\toprule
Policy & Goop & WaterDrop & Sand\\
\midrule
Base & \shortstack{---\\$U:0/1,\ N:0/0$} & \shortstack{$-0.005\pm0.008$\\$U:0/0,\ N:0/0$} & \shortstack{$-0.00004\pm0.003$\\$U:0/0,\ N:0/0$}\\[2pt]
Dense & \shortstack{---\\$U:0/1,\ N:0/0$} & \shortstack{$-0.012\pm0.004$\\$U:0/0,\ N:0/0$} & \shortstack{$-0.003\pm0.003$\\$U:0/0,\ N:0/0$}\\[2pt]
Random25 & \shortstack{$-0.172\pm0.163$\\$U:0/0,\ N:0/0$} & \shortstack{$-0.006\pm0.006$\\$U:0/0,\ N:0/0$} & \shortstack{$-0.011\pm0.019$\\$U:0/0,\ N:0/0$}\\[2pt]
Speed25 & \shortstack{$-0.160\pm0.137$\\$U:0/0,\ N:0/0$} & \shortstack{$-0.005\pm0.006$\\$U:0/0,\ N:0/0$} & \shortstack{$+0.010\pm0.019$\\$U:0/0,\ N:0/0$}\\[2pt]
Cached risk25 & \shortstack{---\\$U:0/1,\ N:0/0$} & \shortstack{$-0.003\pm0.001$\\$U:0/0,\ N:0/0$} & \shortstack{$+0.009\pm0.012$\\$U:0/0,\ N:0/0$}\\[2pt]
RMS25 & \shortstack{$-0.208\pm0.204$\\$U:0/0,\ N:0/0$} & NP & \shortstack{$-0.003\pm0.005$\\$U:0/0,\ N:0/0$}\\[2pt]
\bottomrule
\end{tabular}
\end{table}

Among complete mixed-training policies, speed25 and relative-velocity-RMS25 have full-horizon MSE $.112\pm.027$ and $.125\pm.032$ and beat random25 in every seed. Accuracy gains still leave substantial geometric error: mixed random25, speed25 and RMS25 place $22.16\%$, $19.55\%$ and $20.53\%$ of particles beyond the metadata box by more than $10^{-6}$, versus $0.0632\%$ for matching truth. Their mean trajectory-maximum excursions are $.401$, $.417$ and $.426$, versus $.000722$ for truth. These measure box excursions, not conservation or certified physical failures. The appendix retains every policy and seed, all failures, and these boundary and cost diagnostics (Appendix~\ref{sec:goop-graph-exposure}). A fixed metadata-selected rollout illustrates the remaining physical mismatch in Figure~\ref{fig:qualitative-goop2d}.

For WaterDrop, all 810 H995 autonomous outcomes complete. Under random25, mix-minus-base MSE is $-.0060\pm.0056$, improving in every seed; dense, speed and cached risk also improve in each seed. Yet the autonomous risk-minus-random interaction is $+.00284\pm.00466$, with mixed seed signs. Exposure can therefore improve placement on observed test histories without a consistent autonomous placement gain. Substantial box excursions and increased recorded rollout times remain (Appendix~\ref{sec:waterdrop-110k-exposure}).

\subsection{Why can residual ranking misallocate?}
The earlier WaterDrop controls isolate residual ranking from action value: at 100k updates, faithful/NLL risk correlates with base error ($.357\pm.060$ / $.411\pm.023$), but weakly with actual sparse benefit ($-.021\pm.049$ / $-.052\pm.013$); random has lower one-step error in every seed. Restoring trained self-messages reduces error without reversing that validation/test ordering. Risk's larger alignment gain is outweighed by its squared prediction change in every full-model seed. Autonomous signs vary, with failures and box excursions under both graph conventions (Appendices~\ref{sec:compact-study}, \ref{sec:actual-action-benefit}, \ref{sec:graph-bridge}, \ref{sec:optional-exposure}, \ref{sec:full-results} and~\ref{sec:native-rollout-followup}). These findings motivate testing exposure separately from placement.

\subsection{What computation does the policy require?}
Previous-observed risk includes an extra scoring pass; autonomous cached risk reuses its previous forward output. Their timings measure different computations. Fixed-order autonomous measurements also follow different geometries, so elapsed time and edge counts do not establish a causal speedup. The appendix retains measured cost alongside every policy's error and failure counts. The studies use only three training seeds, reuse earlier evidence to develop follow-up questions, and remain distinct from the original saved-model comparisons.
